\section{既约 Whittaker 范畴}
\markboth{既约 Whittaker 范畴}{}

从现在起，我们假设 $G$ 的中心（记作 $Z_G$）连通。

为了考虑扩展 Whittaker 范畴，我们需要把上面定义的 $\Whit(G)_{glob}$ 换成它的一个满子范畴，记作 $\Whit(G)\rd_{glob}$。

下面是函数论上的类比。如前所述，$D(\Bun\rt_N)$ 是 $N(K)\backslash G(\A)/G(\mathbb O)$ 上函数空间的类比。我们将引入一个范畴，记作 $D(\Bun^{\mathrm{rat},\mathrm{red}}_N)$，它是
\[
  Z_G(K)N(K)\backslash G(\A)/G(\mathbb O)
\]
上函数空间的几何类比。范畴 $\Whit(G)\rd_{glob}$ 将是 $D(\Bun^{\mathrm{rat},\mathrm{red}}_N)$ 中施加了 Whittaker 型条件所得的满子范畴。

\ob{6.1}{blk:6.1}\textbf{范畴。}令 $\overline{\Bun}\rd_N$ 记分类四元组 $(P_G,P_T,\kappa,\gamma\rd)$ 的叠，其中 $(P_G,P_T,\kappa)$ 是 $\Bun_B$ 的一个点，而 $\gamma\rd$ 是下述资料：对 $G$ 的每个单根 $\Ch\alpha_i$，给定线丛的一个非零映射
\[
  \gamma_i\rd : \Ch\alpha_i(P_T)\to\omega_X .
\]

上述叠带有一个群胚的作用，该群胚记作 $\mathcal H_{G,N}$，它实际上是从 $\Bun_B$ 上的群胚 $\mathcal H_{G,B}$ 拉回而来的：只要存在曲线一般点处定义的同构 $P_T\simeq P_T^1$，使 $\kappa$ 与 $\kappa^1$、$\gamma\rd$ 与 $\gamma^{\mathrm{red},1}$ 相配，它就把两点 $(P_G,P_T,\kappa,\gamma)$ 与 $(P_G,P_T^1,\kappa^1,\gamma^1)$ 等同起来。

我们令 $D(\Bun^{\mathrm{rat},\mathrm{red}}_N)$ 记所得的范畴 $D(\overline{\Bun}\rd_N)^{\mathcal H_{G,N}}$。

与 $D(\Bun\rt_B)$ 的情形一样，遗忘函子
\[
  \oblv\rd_N : D(\Bun^{\mathrm{rat},\mathrm{red}}_N)\to D(\overline{\Bun}\rd_N)
\]
有伴随，记作 $\ind\rd_N$，\footnote{原文作「右伴随」。但紧接着的单子陈述 $\oblv\rd_N\circ\ind\rd_N$ 要求 $\ind\rd_N$ 是左伴随（与 §1.2 的 $\oblv_B\dashv\ind_B$ 一致）；故此处「右」字应为笔误。译文照录原文。}而单子
\[
  \oblv\rd_N\circ\ind\rd_N : D(\overline{\Bun}\rd_N)\to D(\overline{\Bun}\rd_N)
\]
由关于 $\mathcal H_{G,N}$ 取平均给出，即函子 $\Av_{\mathcal H_{G,H}}$。\footnote{原文此处的下标是 $\mathcal H_{G,H}$，全书他处并无 $\mathcal H_{G,H}$；由上下文，应为 $\mathcal H_{G,N}$。译文照录原文。}

\ob{6.2}{blk:6.2}\textbf{Whittaker 条件。}略微滥用记号，我们把从 $\Bun_B$ 拉回到 $\overline{\Bun}\rd_N$ 的同名群胚仍记作 $\mathcal H_N$。注意，$\gamma\rd$ 的资料现在定义出典范特征
\[
  \chi : \mathcal H_N\to\Ga .
\]
由于 $\mathcal H_N$ 与 $\mathcal H_{G,N}$ 的作用交换，我们可以考虑范畴
\[
  \Whit(G)\rd_{glob} :=
  D(\overline{\Bun}\rd_N)^{(\mathcal H_N,\chi),\mathcal H_{G,N}}
  = D(\Bun^{\mathrm{rat},\mathrm{red}}_N)^{(\mathcal H_N,\chi)},
\]
它是 $D(\overline{\Bun}\rd_N)^{\mathcal H_{G,N}}=D(\Bun^{\mathrm{rat},\mathrm{red}}_N)$ 的满子范畴。

我们有伴随函子对
\[
  \ind_{\Whit\rd,glob} : D(\Bun^{\mathrm{rat},\mathrm{red}}_N)
  \rightleftarrows \Whit(G)\rd_{glob} : \oblv_{\Whit\rd,glob}.
\]

\ob{6.2.1}{blk:6.2.1}\emph{典范对象。}范畴 $\Whit(G)\rd_{glob}$ 有典范的单位对象，记作 $\one_{\Whit(G)\rd_{glob}}$：它由把 $\ind\rd_N$ 作用于 $D(\overline{\Bun}\rd_N)$ 中某个典范的 $(\mathcal H_N,\chi)$-等变对象而得。

所说的对象是如下特征层的 $!$-（同时也是 $\ast$-）延拓：该特征层定义在 $\Bun\rd_N\subset\overline{\Bun}\rd_N$ 上，其中 $\Bun\rd_N$ 是 $\overline{\Bun}\rd_N$ 中 Plücker 数据 $\kappa$ 既无零点也无极点、且 $\gamma_i\rd$ 均为同构的那个轨迹。特征层由沿留数映射 $\Bun\rd_N\to\Ga$ 拉回 Artin-Schreier 层而得。

\ob{6.3}{blk:6.3}\textbf{$\Whit(G)_{glob}$ 与 $\Whit(G)\rd_{glob}$ 的关系。}

在 [GL:Ran] 中我们将证明：

\begin{othm}{命题}{6.3.1}{prop:6.3.1}
范畴 $\Whit(G)_{glob}$ 带有（刚性）幺半范畴 $D(\Gr_{Z_G,\mathrm{Ran}_X,un})$ 的自然作用，且有典范等价
\[
  \Whit(G)\rd_{glob}\simeq
  \Whit(G)_{glob}\underset{D(\Gr_{Z_G,\mathrm{Ran}_X,un})}{\otimes} D(\Bun_{Z_G}).
\]
\end{othm}

特别地，有伴随函子对
\[
  \Whit(G)_{glob}\rightleftarrows\Whit(G)\rd_{glob},
\]
其右伴随全忠实。

注意，$\one_{\Whit(G)_{glob}}$ 在上述函子 $\Whit(G)_{glob}\to\Whit(G)\rd_{glob}$ 下的像，正是上面那个对象 $\one_{\Whit(G)\rd_{glob}}$。

\ob{6.3.2}{blk:6.3.2}上述命题与定理 5.3.3 结合，给出 $\Whit(G)\rd_{glob}$ 的如下谱描述：

\begin{othm}{推论}{6.3.3}{cor:6.3.3}
我们有
\[
  \Whit(G)\rd_{glob}\simeq
  \Rep(\Ch G,\mathrm{Ran}_X,un)
  \underset{\Rep(\Ch G/[\Ch G,\Ch G],\mathrm{Ran}_X,un)}{\otimes}
  \QCoh\bigl(\LocSys_{\Ch G/[\Ch G,\Ch G]}\bigr).
\]
\end{othm}

\begin{annote}
命题 6.3.1 与推论 6.3.3 是同一件事的两侧。几何侧说：既约 Whittaker 范畴是把 $\Whit(G)_{glob}$ 沿 $\Gr_{Z_G}$ 的（Hecke）作用「张量缩小」到 $D(\Bun_{Z_G})$；谱侧说：这相当于把 $\Rep(\Ch G)$ 换成它商去换位子群 $\Ch G/[\Ch G,\Ch G]$ 后的表示，并把 $\QCoh(\LocSys_{\Ch G})$ 换成 $\Ch G/[\Ch G,\Ch G]$ 的局部系统叠上的拟凝聚层。两边都体现「$Z_G$ 连通」这一假设的作用：它保证中心部分可以在两侧一致地剥离。\par\vspace{0.15em}{\raggedleft 【来源：原书 p.11 §6.3--6.3.2（A）】\par}
\end{annote}
